\documentclass[11pt]{article}
\usepackage[margin=1in]{geometry}
\usepackage{amsmath,amssymb,amsthm,hyperref,enumitem}
\newtheorem{theorem}{Theorem}
\newtheorem{lemma}{Lemma}
\newcommand{\e}{e}
\title{An alternating Tur\'an inequality for the Euler numbers}
\author{Principia Math In Progress}
\date{21 September 2026; source audit updated 22 September 2026}
\begin{document}
\maketitle
\noindent\textbf{MathDB 338918.} Full proof of the assertion in Conjecture 3.3
of Sokal's accepted manuscript \cite{Sokal}. The source audit uses the complete
author accepted version deposited at UCL; the separately typeset journal
version remains unread. Review is by the authoring agent only.

\section*{Task 1: the full inequality}
Define the Euler zigzag numbers by
\[
 \sec t+\tan t=\sum_{n\ge0}E_n\frac{t^n}{n!},\qquad e_n=E_n/n!.
\]
Thus $E_0,E_1,E_2,E_3,E_4=1,1,1,2,5$; these are the positive
secant and tangent numbers interleaved.
\begin{theorem}\label{main}
For all integers $n\ge0$ and $j,k\ge1$,
\begin{equation}\label{target}
 (-1)^{n-1}\bigl(e_ne_{n+j+k}-e_{n+j}e_{n+k}\bigr)>0.
\end{equation}
\end{theorem}

Put $c=4/\pi$, $a=2/\pi$, and define
\[
 B(s)=\sum_{r\ge0}\frac{(-1)^r}{(2r+1)^s}\quad(s\ge1),\qquad
 L(s)=\sum_{r\ge0}\frac1{(2r+1)^s}\quad(s>1).
\]
The classical partial fraction expansions of secant and tangent give
\begin{equation}\label{parity}
 e_n=ca^n\begin{cases}B(n+1),&n\text{ even},\\L(n+1),&n\text{ odd}.
 \end{cases}
\end{equation}
These formulas, including the conditionally convergent case $B(1)=\pi/4$,
are recorded in \cite[eqs. (3.3)--(3.5), (3.16)]{Sokal}.

\begin{lemma}\label{mono}
For $s\ge1$, $0<B(s)<1$ and $B(s+2)>B(s)$.
For $s>1$, $L(s)>1$ and $L(s+2)<L(s)$.
\end{lemma}
\begin{proof}
The bounds for $B$ follow from the alternating series, and those for $L$
from its positive terms. The decrease of $L$ is termwise, strictly at $r=1$.
For $B$ subtract convergent partial sums to obtain
\[
 B(s+2)-B(s)=\sum_{r\ge1}(-1)^{r+1}h_s(2r+1),\qquad
 h_s(x)=x^{-s}-x^{-s-2}.
\]
For $x\ge3$, $h_s(x)>0$, $h_s(x)\to0$, and
\[
 h_s'(x)=x^{-s-3}\bigl(s+2-sx^2\bigr)<0.
\]
Pairing consecutive terms proves that this alternating sum is strictly
positive. No differentiation of an infinite series is used at $s=1$.
\end{proof}

We recall with a proof the two known same-parity properties used below.
They are established in \cite[Section 3, Proposition 3.2]{Sokal}.
\begin{lemma}\label{same}
The positive sequence $(e_{2m})_{m\ge0}$ is strictly log-concave;
$(e_{2m+1})_{m\ge0}$ is strictly log-convex.
Consequently, for integers $m\ge0$ and $u,v\ge1$,
\[
 e_{2m}e_{2(m+u+v)}<e_{2(m+u)}e_{2(m+v)},\qquad
 e_{2m+1}e_{2(m+u+v)+1}>e_{2(m+u)+1}e_{2(m+v)+1}.
\]
\end{lemma}
\begin{proof}
Write $x_r=4/[\pi^2(2r+1)^2]$. The cosine product yields
\[
 \sum_{m\ge0}e_{2m}z^m=\prod_{r\ge0}(1-x_rz)^{-1}.
\]
Here is an elementary strictness argument. For a sequence $h$ set
$T(h)_{ij}=h_{j-i}$, with $h_t=0$ for $t<0$.
Every two-by-two minor of $T(1,x,x^2,\ldots)$ is nonnegative: if all
four entries are nonzero its determinant is zero, and in the other cases
it is either zero or the product of its diagonal entries.
Convolution multiplies these upper triangular matrices. The two-by-two
Cauchy--Binet formula therefore preserves nonnegative minors; each sum
here is finite.

For the first two factors the coefficient is
$h_m=\sum_{i=0}^m x_0^i x_1^{m-i}$, and direct algebra gives
\[
 h_m^2-h_{m-1}h_{m+1}=(x_0x_1)^m>0\qquad(m\ge1).
\]
Multiply their Toeplitz matrix on the left by the matrices of the
remaining factors. For rows $0,1$ and columns $m,m+1$, the intermediate
indices $0,1$ in Cauchy--Binet contribute precisely the displayed positive
minor: the left matrix has leading minor $1$. All other contributions
are nonnegative. The finite products thus have this same positive lower
bound. Since $\sum x_r<\infty$, their coefficients converge to those of
the infinite product, proving strict log-concavity of $e_{2m}$.

The tangent expansion gives
$e_{2m+1}=2\sum_{r\ge0}x_r^{m+1}$.
Strict Cauchy--Schwarz for this positive measure, supported at distinct
positive points, proves strict log-convexity. All sums converge.
Finally, strict log-concavity means that consecutive ratios strictly
decrease. Multiplying $u$ such ratios, starting at $m$ and at $m+v$,
proves the first rectangular inequality. Increasing ratios prove the
second one.
\end{proof}

\begin{proof}[Proof of Theorem \ref{main}]
In every product in \eqref{target}, the positive factor $c^2a^{2n+j+k}$
from \eqref{parity} is the same, so it can be canceled.

If $j,k$ are both even, apply Lemma \ref{same} with $u=j/2$, $v=k/2$.
If $j,k$ are both odd and $n$ is even, the first product after cancellation
is a product of two $B$ values, hence is less than $1$; the second is a
product of two $L$ values, hence is greater than $1$. The difference is
negative. For odd $n$ the two roles interchange, so it is positive.

In the remaining case, symmetry allows us to take $j$ even and $k$ odd.
If $n$ is even, the difference after cancellation is
\[
 B(n+1)L(n+j+k+1)-B(n+j+1)L(n+k+1)<0.
\]
Indeed, both factors of its first product are strictly smaller than
the corresponding factors of its second, by Lemma \ref{mono} iterated
$j/2$ times. If $n$ is odd, the difference is
\[
 L(n+1)B(n+j+k+1)-L(n+j+1)B(n+k+1)>0,
\]
because both factors of its first product are strictly larger. All $L$
arguments used are at least $2$, all $B$ arguments at least $1$.
This proves every case, including $n=0$.
\end{proof}

\paragraph{Attribution and verification.}
The partial fraction expansions and the two same-parity properties are
prior results, explicitly attributed above. The contribution is the
complete parity argument for the conjectured three-parameter inequality.
The accompanying exact verifier checks Euler coefficients by independent
recurrences, finite instances of all parity classes, the product-minor
identity, and the alternating-series comparison. Finite checks supplement
the proof; they do not prove its unbounded quantifiers.

\section*{Four proposed follow-up tasks}
These directions are not claimed as solved.
\begin{enumerate}[label=Task \arabic*:,start=2,leftmargin=4.5em]
\item For $b_n(q)=\sum_{r\ge0}(-1)^{r(n+1)}(q+2r)^{-n-1}$, determine
all $q>0$ for which the analogue of \eqref{target} holds for every $n,j,k$.
\item Determine the signs of the consecutive Hankel determinants
$\det(e_{n+i+j})_{0\le i,j<r}$ for every $n\ge0$ and $r\ge3$.
\item For fixed $j,k$, find the exact leading asymptotic of
$(-1)^{n-1}(e_ne_{n+j+k}-e_{n+j}e_{n+k})$ separately on even and odd $n$,
with an explicit error bound.
\item Construct a combinatorial proof of \eqref{target} after multiplication
by $n!(n+j)!(n+k)!(n+j+k)!$, using alternating permutations.
\end{enumerate}

\begin{thebibliography}{1}
\bibitem{Sokal} A. D. Sokal, \emph{The Euler and Springer numbers as moment
sequences}, \href{https://arxiv.org/abs/1804.04498}{arXiv:1804.04498v3}
(2018), Conjecture 3.3, p. 14. Published in \emph{Expositiones Mathematicae}
\textbf{38} (2020), 1--26,
\href{https://doi.org/10.1016/j.exmath.2018.08.001}{doi:10.1016/j.exmath.2018.08.001}.
All 30 pages of the current manuscript were read. The UCL deposit explicitly
identifies the same file as the accepted version:
\url{https://discovery.ucl.ac.uk/id/eprint/10059424/}.
The separately typeset journal version remains unread.
\end{thebibliography}
\end{document}
